\chapter{Formal Definitions}\label{app:definitions}

This appendix collects the formal definitions on which the design (\cref{ch:design}) and its
analysis rest. Each is stated in the notation of the Notation table, and each records where it is
used in the thesis. Definitions of standard primitives are given as they appear in the specifications
they are drawn from and are cited accordingly; the dual-index scheme of \cref{sec:def-dualindex} is
the one construction original to this work.

% ============================================================================
\section{Cryptographic primitives}\label{sec:def-primitives}

\begin{definition}[Collision-resistant hash function]\label{def:hash}
\intuition{A function whose outputs are short fingerprints that no efficient adversary can make
collide.}
A hash function $\Hash:\{0,1\}^*\to\{0,1\}^{256}$ is collision-resistant if for every probabilistic
polynomial-time adversary $\Adv$ the probability that $\Adv$ outputs $x\neq x'$ with
$\Hash(x)=\Hash(x')$ is negligible in the security parameter $\secpar$. The instantiation used
throughout is SHA-256. Used in every authenticated structure of \cref{sec:def-structures}.
\end{definition}

\begin{definition}[Key-derivation function, HKDF]\label{def:hkdf}
\intuition{A two-step ``extract then expand'' procedure that turns a high-entropy secret into any
number of independent keys.}
Following RFC~5869~\cite{rfc5869}, $\HKDFex(\mathit{salt}, \mathit{IKM})\to \mathit{PRK}$ extracts a
pseudorandom key from input keying material, and $\HKDFexp(\mathit{PRK}, \mathit{info}, L)\to
\mathit{OKM}$ expands it into $L$ bytes of output bound to a context string $\mathit{info}$. Both are
instantiated with HMAC-SHA256. Used in the dual-index label derivation
(\cref{def:dualindex-derive}).
\end{definition}

\begin{definition}[Digital signature, EUF-CMA]\label{def:sig}
\intuition{A signature scheme that no adversary can forge, even after seeing signatures of its
choice.}
A signature scheme $(\mathsf{KeyGen},\mathsf{Sign},\mathsf{Verify})$ is existentially unforgeable
under chosen-message attack if no probabilistic polynomial-time adversary, given the public key and a
signing oracle, can produce a valid signature on a message it did not query, except with negligible
probability in $\secpar$. The instantiation is Ed25519 (RFC~8032~\cite{rfc8032}), used to sign the
signed tree head (\cref{def:sth}).
\end{definition}

% ============================================================================
\section{Authenticated data structures}\label{sec:def-structures}

\begin{definition}[Merkle log tree and tree hash]\label{def:merkle-log}
\intuition{A tamper-evident list: one short root fingerprint commits to every entry in order, and any
change to the history changes the root.}
For an ordered list of entries $D=(d_0,\dots,d_{n-1})$, the Merkle Tree Hash
$\MTH(D)$ is defined recursively as in RFC~6962~\cite{rfc6962}:
$\MTH(\varnothing)=\Hash(\text{``''})$, $\MTH(d_0)=\Hash(\texttt{0x00}\concat d_0)$, and for $n>1$,
\[
  \MTH(D)=\Hash\!\big(\texttt{0x01}\concat \MTH(D[0{:}k]) \concat \MTH(D[k{:}n])\big),
\]
where $k$ is the largest power of two strictly less than $n$. The domain-separation prefixes
\texttt{0x00} (leaf) and \texttt{0x01} (internal node) prevent second-preimage attacks across the two
node types. The log root at epoch $t$ is $\logroot{t}=\MTH(D_t)$. Used in
\cref{ssec:design-structures}.
\end{definition}

\begin{definition}[Inclusion proof]\label{def:inclusion}
\intuition{The list of sibling hashes on the path from a leaf to the root, from which anyone can
recompute the root and confirm the leaf is present.}
An inclusion proof $\pincl$ for the entry at index $i$ in a log of size $n$ is the audit path of
RFC~9162~\cite{rfc9162}: the sequence of subtree hashes that, combined with $d_i$ along the path from
leaf $i$ to the root, recomputes $\MTH(D)$. Verification recomputes the candidate root from $d_i$ and
$\pincl$ and accepts iff it equals the committed $\logroot{t}$. Used in \cref{ssec:design-verify}.
\end{definition}

\begin{definition}[Consistency proof]\label{def:consistency}
\intuition{A short proof that a later version of the log is an append-only extension of an earlier
one --- nothing was reordered, modified, or deleted.}
For sizes $m\le n$, a consistency proof $\pcons(m,n)$ is the set of subtree hashes of
RFC~9162~\cite{rfc9162} that proves the size-$m$ tree is a prefix of the size-$n$ tree, i.e.\ that
$\MTH(D[0{:}m])$ is consistent with $\MTH(D[0{:}n])$. Verification accepts iff the two committed roots
can be reconstructed from the proof. Used in \cref{ssec:design-catchup}.
\end{definition}

\begin{definition}[Sparse Merkle prefix tree]\label{def:prefix-tree}
\intuition{A Merkle tree over the whole space of possible labels, so that a label's slot has a fixed
position and even an \emph{absent} label has a provable (empty) location.}
A sparse Merkle prefix tree is a binary Merkle tree of depth $256$ indexed by 256-bit labels
$\ell$, in which the leaf at position $\ell$ commits to $\Hash(\ell\concat \mathit{value})$ and every
empty subtree collapses to a fixed default hash. Following the IETF Key Transparency
draft~\cite{ietf-keytrans-protocol} and CONIKS~\cite{melara2015coniks}, the tree is path-compressed
(single-child chains are elided) and proofs are transmitted with a bitmap recording which siblings
are non-default, so a proof carries only the siblings that actually exist. The prefix root at epoch
$t$ is $\prefroot{t}$. Used in \cref{ssec:design-structures}.
\end{definition}

\begin{definition}[Inclusion and non-inclusion proof in the prefix tree]\label{def:noninclusion}
\intuition{Because every label has a fixed position, one can prove a label is \emph{absent} by
exhibiting the (default) contents of its slot, just as one proves presence.}
A prefix-tree proof for label $\ell$ is the co-path of sibling hashes from the leaf at position
$\ell$ to the prefix root $\prefroot{t}$. An \emph{inclusion} proof recomputes $\prefroot{t}$ from
$\Hash(\ell\concat \mathit{value})$ and the co-path; a \emph{non-inclusion} proof recomputes it from
the default (empty) leaf value at position $\ell$, establishing that no value is bound at $\ell$
\cite{melara2015coniks,ietf-keytrans-protocol}. Proven non-inclusion of the warning slot is what
makes warning censorship detectable (\cref{def:dualindex-check}). Used in \cref{ssec:design-verify}.
\end{definition}

\begin{definition}[Signed tree head]\label{def:sth}
\intuition{A single signed value that publicly commits the operator to one view of the whole
directory and its history at a given epoch.}
The signed tree head at epoch $t$ is
$\STH_t = \big(\,t,\ \mathit{size}_t,\ \logroot{t},\ \mathit{ts}_t,\ \Hash(\mathit{config})\,\big)$
together with an Ed25519 signature over that tuple under the operator's signing key
(\cref{def:sig}), in the manner of the Certificate Transparency signed tree head~\cite{rfc6962}. A
client or auditor recomputes the signed message from the head it received and checks the signature
before trusting any proof anchored to it. Used in \cref{ssec:design-verify,ssec:design-epoch}.
\end{definition}

% ============================================================================
\section{Verifiable random function}\label{sec:def-vrf}

\begin{definition}[Verifiable random function]\label{def:vrf}
\intuition{A keyed function whose outputs look random to everyone but the key holder, yet come with a
proof that the output was computed correctly.}
A verifiable random function is a triple $(\VRFkg,\VRFeval,\VRFvfy)$~\cite{rfc9381}, comprising
\[
\begin{aligned}
  \VRFkg(1^\secpar) &\to (\mathit{pk},\mathit{sk}), \\
  \VRFeval(\mathit{sk},\alpha) &\to (\beta,\pvrf) \quad\text{(an output $\beta$ and a proof $\pvrf$)}, \\
  \VRFvfy(\mathit{pk},\alpha,\pvrf) &\to \beta \ \text{ or }\ \bot.
\end{aligned}
\]
The instantiation is ECVRF-EDWARDS25519-SHA512-TAI, suite \texttt{0x03}, of
RFC~9381~\cite{rfc9381}. Used to derive global directory labels in \cref{ssec:design-structures}.
\end{definition}

\begin{definition}[VRF uniqueness and pseudorandomness]\label{def:vrf-security}
\intuition{Uniqueness: a given key and input admit only one verifiable output. Pseudorandomness:
without the secret key, that output is indistinguishable from random.}
For \emph{full uniqueness}, for every (even adversarially chosen) public key $\mathit{pk}$ and input
$\alpha$ there is at most one $\beta$ for which some proof verifies. For \emph{pseudorandomness}, no
probabilistic polynomial-time adversary wins the experiment $\Expt{vrf\text{-}prand}{\Adv}$ of
RFC~9381~\cite{rfc9381} with advantage non-negligible in $\secpar$: the adversary, given
$\mathit{pk}$ and a $\VRFeval$ oracle, must distinguish the true output $\beta^\star=\VRFeval(\mathit{sk},\alpha^\star)$
on a fresh challenge $\alpha^\star$ from a uniform string. Uniqueness underlies label integrity
in the verification chain (goal G1 of \cref{ch:threat}); pseudorandomness underlies directory
privacy (goal G3).
\end{definition}

% ============================================================================
\section{Epochs and windowed consistency}\label{sec:def-epoch}

\begin{definition}[Epoch]\label{def:epoch}
\intuition{A committed snapshot of the directory; between snapshots the live tree is mutated, and a
commit freezes one version into the public history.}
An epoch is a committed version of the directory. An
$\mathsf{update}(\ell,\mathit{value})$ mutates the live prefix tree in place; a
$\mathsf{commit}$ closes the epoch by appending the current prefix root $\prefroot{t}$ as one leaf
of the log (\cref{def:merkle-log}) and publishing the signed tree head $\STH_t$ (\cref{def:sth}).
Lookups are always served against the last committed epoch, never the live tree. Used in
\cref{ssec:design-epoch}.
\end{definition}

\begin{definition}[Windowed consistency proof]\label{def:windowed}
\intuition{A single consistency proof that lets a long-absent client re-synchronise in one step,
rather than one proof per missed epoch.}
For a client that last verified the log at size $m$ and returns when the log has size $n$, spanning
$W$ intervening epochs, the windowed consistency proof is the single proof $\pcons(m,n)$ of
\cref{def:consistency}. Used in \cref{ssec:design-catchup}.
\end{definition}

\begin{lemma}[Windowed proof size]\label{lem:windowed-size}
The windowed consistency proof $\pcons(m,n)$ consists of at most $\lceil\lg n\rceil + \lceil\lg
m\rceil$ hash values --- logarithmic in the endpoint sizes rather than linear in the number of
intervening epochs $W$.
\end{lemma}
\begin{proof}[Proof sketch]
By the structure of the RFC~9162 consistency proof~\cite{rfc9162}, the proof comprises the subtree
hashes needed to reconstruct both $\MTH(D[0{:}m])$ and $\MTH(D[0{:}n])$; the number of such nodes is
bounded by the depths of the two trees, $\lceil\lg m\rceil$ and $\lceil\lg n\rceil$. The bound
depends only on the endpoint sizes $m$ and $n$, growing logarithmically with them rather than
linearly with the $W$ epochs between the two visits, which is the property exploited in
\cref{ssec:design-catchup}.
\end{proof}

% ============================================================================
\section{Signal protocol primitives}\label{sec:def-signal}

\begin{definition}[X3DH shared secret]\label{def:x3dh}
\intuition{An asynchronous handshake that lets Alice derive a shared secret with an offline Bob from
his published pre-keys.}
The Extended Triple Diffie--Hellman handshake~\cite{x3dh} derives a shared secret $\sharedS$ from a
combination of Diffie--Hellman values over the parties' identity keys $\pkid{\cdot}$, Bob's signed
pre-key $\spk{\B}$, an optional one-time pre-key $\opk{\B}{i}$, and Alice's ephemeral key
$\ek{\A}$; the post-quantum variant PQXDH~\cite{pqxdh} augments this with a KEM. In this thesis
$\sharedS$ is the entropy source for the dual-index labels (\cref{def:dualindex-derive}). Introduced
in \cref{ssec:design-labels}.
\end{definition}

\begin{definition}[Double Ratchet]\label{def:doubleratchet}
\intuition{A chain of key-derivation steps that gives every message a fresh key and heals the session
after a compromise.}
The Double Ratchet~\cite{doubleratchet} maintains a root chain and per-direction sending and
receiving chains, deriving a fresh message key per message and mixing in new Diffie--Hellman outputs
as the conversation proceeds. Its per-epoch ratchet state is the input to the continuous operating
mode (\cref{def:cka}). Referenced in \cref{ssec:design-modes}.
\end{definition}

\begin{definition}[Continuous key agreement and continuous authentication]\label{def:cka}
\intuition{An abstraction of a ratcheting protocol as a stream of evolving shared keys, together with
a notion of authenticating that stream continuously.}
The Authenticated Continuous Key Agreement framework~\cite{dowling2023acka} models a ratcheting
protocol as a continuous key agreement $\CKA$ producing a sequence of per-epoch keys, and defines a
continuous entity authentication (CEA) notion, with security game $\CEAsec$, that captures detection of an active adversary who
injects state mid-session. The framework leaves its append-only log component $\Logsec$ abstract; this
thesis instantiates that component with the structures of \cref{sec:def-structures} and realises
continuous authentication through the continuous mode of \cref{ssec:design-modes}. The dual-index
warning scheme below is an \emph{extension} of this idea --- it authenticates by directory
\emph{location} rather than by chained fingerprint value --- and is not an instance of $\CEAsec$.
\end{definition}

% ============================================================================
\section{The dual-index warning scheme}\label{sec:def-dualindex}

The following definitions specify the construction original to this work
(\cref{ssec:design-labels}). They rely on the primitives above and introduce no new assumptions.

\begin{definition}[Dual-index derivation]\label{def:dualindex-derive}
\intuition{From one shared secret, two parties independently derive the same three directory slots:
one lookup slot each and one shared warning slot.}
$\mathsf{DeriveIndices}(\sharedS, \mathit{id}_{\mathsf{self}}, \mathit{id}_{\mathsf{peer}})$ first
extracts $\mathit{PRK}=\HKDFex(\text{salt}=\text{\texttt{"duet-dual-index"}},\ \sharedS)$ and then returns the three 256-bit labels
\begin{align*}
  \lblself &= \HKDFexp(\mathit{PRK},\ \texttt{"lookup"}\concat \mathit{id}_{\mathsf{self}}\concat \mathit{id}_{\mathsf{peer}},\ 32),\\
  \lblpeer &= \HKDFexp(\mathit{PRK},\ \texttt{"lookup"}\concat \mathit{id}_{\mathsf{peer}}\concat \mathit{id}_{\mathsf{self}},\ 32),\\
  \lblwarn &= \HKDFexp(\mathit{PRK},\ \texttt{"warning"}\concat \min(\mathit{id}_\A,\mathit{id}_\B)\concat \max(\mathit{id}_\A,\mathit{id}_\B),\ 32),
\end{align*}
using HKDF-SHA256 (\cref{def:hkdf}). The canonical ordering of identifiers in $\lblwarn$ makes the
warning label symmetric, so both parties address the same warning slot, while the lookup labels are
deliberately asymmetric. Specified in \cref{ssec:design-labels}.
\end{definition}

\begin{definition}[Register and write-warning]\label{def:dualindex-write}
\intuition{A party publishes its key at its own lookup slot, and raises the alarm by writing to the
shared warning slot.}
$\mathsf{Register}(\lblself,\mathit{pk})$ binds the party's public key at its own lookup label in the
directory. $\mathsf{WriteWarning}(\lblwarn)$ writes a non-empty value to the shared warning slot.
Both take effect at the next committed epoch (\cref{def:epoch}). Specified in
\cref{ssec:design-detection}.
\end{definition}

\begin{definition}[Check-warning and the verification verdict]\label{def:dualindex-check}
\intuition{A client accepts a peer key only if the peer's lookup slot verifies against a signed head
\emph{and} the shared warning slot is provably empty.}
$\mathsf{CheckWarning}(\lblwarn,\STH_t)$ returns \textsc{raised} if an inclusion proof binds a
non-empty value at $\lblwarn$ under $\STH_t$, and \textsc{clear} if a non-inclusion proof
(\cref{def:noninclusion}) establishes the slot is empty; a response with neither valid proof is
rejected. A peer key is accepted only when the peer-lookup proof verifies through the chain of
\cref{ssec:design-verify} and $\mathsf{CheckWarning}$ returns \textsc{clear}. Specified in
\cref{ssec:design-verify}.
\end{definition}

\begin{definition}[Dual-index detection experiment]\label{def:dualindex-detect}
\intuition{The adversary wins only if it makes a client accept a substituted key while the warning
slot stays provably empty --- i.e.\ if it substitutes without tripping the wire.}
Let $\Adv$ be a probabilistic polynomial-time adversary controlling the server and the network under
the model of \cref{ch:threat}. In the experiment $\Expt{dual\text{-}det}{\Adv}$, $\Adv$ interacts with
an honest client and peer and wins if it causes the client to accept a peer key
$\mathit{pk}'\neq\mathit{pk}_{\B}$ under a signed head $\STH_t$ for which $\mathsf{CheckWarning}$
returns \textsc{clear}. The scheme achieves key-substitution detection (goal G1 of \cref{ch:threat})
if $\advantage{dual\text{-}det}{\Adv}=\Pr[\Expt{dual\text{-}det}{\Adv}=1]$ is negligible in
$\secpar$, under the collision resistance of $\Hash$ (\cref{def:hash}), the unforgeability of the
signature (\cref{def:sig}), and the uniqueness of the VRF (\cref{def:vrf-security}). The
narrative argument for this claim is given in \cref{ssec:design-detection}; a full reduction is left
to future work (\cref{ch:discussion}).
\end{definition}